% Requires graphicx and amsmath. Put these PDFs beside the manuscript or set
% \graphicspath to this directory. Final figure numbering is assigned by LaTeX.
% Figure A below is former B2; Figure B below is former B1.
% Displayed t = source task + 1 = source T. Source IDs remain unchanged.
\begin{figure*}[t]
\centering
\includegraphics[width=\textwidth]{Figure_A_realized_future.pdf}
\caption{Realized future-learning diagnostics on ImageNet-A for EWCLoRA and FS-EWCLoRA under the formal EWC configuration $\lambda=2000$.
From left to right: effective-weight drift, Task-1 signed test-loss change, Task-1 test accuracy, and path Fisher sensitivity.
All horizontal axes show the number of learned tasks $t=1,\ldots,10$.
For each method and seed, $\Delta\mathbf w_t=\mathbf w_t-\mathbf w_1$ is the displacement from its own Task-1 anchor in the 24 effective Q/V matrices and $d_t=\|\Delta\mathbf w_t\|_2$.
The signed loss change is $D_{1,t}^{\mathrm{test}}=\widehat{\mathcal L}_{1,\mathrm{test}}(\mathbf w_t)-\widehat{\mathcal L}_{1,\mathrm{test}}(\mathbf w_1)$; negative values are retained.
Loss and accuracy $A_{1,t}^{\mathrm{test}}$ use all 307 Task-1 test examples and the fixed 20-class Task-1 head.
The first three panels contain all ten endpoints per method and seed (60 endpoint records).
Path sensitivity is $\overline S_{F,1}^{\mathrm{path}}(t)=\int_0^1\mathbf v_t^\top F_{1,\mathrm{tr}}^{\mathrm{emp}}(\mathbf w_1+s\Delta\mathbf w_t)\mathbf v_t\,\mathrm ds$, where $\mathbf v_t=\Delta\mathbf w_t/d_t$.
It uses raw true-label per-example empirical Fisher on all 1,259 Task-1 training examples, without EWC coefficient scaling, trace normalization, damping, floor, or added regularization.
Straight anchor-to-endpoint chords are independently measured at $t=2,\ldots,10$; no estimate is plotted at $t=1$ because $d_1=0$.
Adaptive nested Simpson quadrature uses 9 or 17 measured nodes per chord, for 54 chords and 502 nodes in total.
The chords are not the optimizer's chronological trajectory, and their directional sensitivity is not an endpoint Fisher trace.
Curves show arithmetic means over training seeds 0, 1, and 2, with each seed's integral divided by its own squared displacement before averaging.
Standard-deviation shading is omitted; sample standard deviations remain available in the source CSV.}
\label{fig:rq2-realized-future}
\end{figure*}

\begin{figure*}[t]
\centering
\includegraphics[width=\textwidth]{Figure_B_controlled_neighborhood.pdf}
\caption{Controlled neighborhood diagnostics on ImageNet-A around the formal seed-0 Task-1 anchors of EWCLoRA and FS-EWCLoRA under the EWC configuration $\lambda=2000$.
From left to right: Task-1 train accuracy, mean positive loss increase, neighborhood Fisher trace, and measured tolerance radius.
The first three horizontal axes show controlled perturbation radius $d$; the fourth shows allowed loss increase $\tau$.
The direction-averaged accuracy is $\overline A_1^{\mathrm{nbr}}(d)=K_D^{-1}\sum_k A_{1,\mathrm{tr}}(\mathbf w_1+d\mathbf u_k)$ in percent.
The positive-loss statistic is $\overline D_1^+(d)=K_D^{-1}\sum_k[\widehat{\mathcal L}_{1,\mathrm{tr}}(\mathbf w_1+d\mathbf u_k)-\widehat{\mathcal L}_{1,\mathrm{tr}}(\mathbf w_1)]_+$: the positive part is taken for each direction before averaging.
Loss and accuracy use all 1,259 Task-1 training examples, the fixed 20-class Task-1 head, $K_D=50$ paired random unit directions, and 107 measured radii per method; the coincident directions at $d=0$ require one anchor evaluation.
The raw neighborhood Fisher statistic is $\overline R_{F,1}^{\mathrm{nbr}}(d)=K_F^{-1}\sum_k\operatorname{tr}F_{1,\mathrm{tr}}^{\mathrm{emp}}(\mathbf w_1+d\mathbf u_k)$, using 256 fixed training examples, $K_F=16$ paired directions, and 42 measured radii per method.
It averages true-label per-example gradient squared norms in the 24 effective Q/V matrices, without EWC coefficient scaling, trace normalization, or damping.
Anchors use rank-10 LoRA; diagnostic perturbations occupy the full effective-weight space without a rank-10 constraint.
The measured tolerance radius is $\widehat d_\tau=\max\{d_j:\overline D_1^+(d_i)\leq\tau\ \text{for all }i\leq j\}$, computed on the measured grid without interpolation or a continuous-radius certificate.
All eleven thresholds through $\tau=0.1$ are retained.
Radii $d=21,22,23,24$ were not measured, limiting tolerance resolution; connecting lines are visual guides.
The measured Fisher peak at $d=17$ is retained.
These are direction means from one training seed, with no cross-seed uncertainty band.}
\label{fig:rq2-controlled-neighborhood}
\end{figure*}
